El análisis y diseño de algoritmos es un área fundamental en Ciencias de la Computación, debido a que permite evaluar la eficiencia de diferentes métodos para resolver problemas de carácter computacional. El estudio de la complejidad temporal y espacial de los algoritmos es esencial para poder  elegir la solución más adecuada teniendo en cuenta los recursos disponibles y el tamaño de los datos de entrada.

Dentro de este contexto, los algoritmos de ordenamiento y de multiplicación de matrices son herramientas clásicas que han sido ampliamente estudiadas, debido a su aplicación en diversos dominios, como bases de datos, gráficos computacionales y aprendizaje automático. Si bien existen soluciones eficientes en teoría, en la práctica pueden variar según las condiciones de entrada y las características específicas de cada implementación.

El propósito de este informe es comparar experimentalmente diferentes algoritmos de ordenamiento (SelectionSort, MergeSort, QuickSort y std::sort) y de multiplicación de matrices (Naive y Strassen), evaluando su desempeño en términos del tiempo de ejecución y el uso de memoria. Este análisis busca identificar patrones de comportamiento y contrastar la eficiencia teórica con el rendimiento experimental, destacando fortalezas y debilidades de cada enfoque.

Este trabajo tiene como objetivo reforzar el entendimiento de los conceptos de complejidad de algoritmos, aplicándolos en un contexto práctico y experimental.



\begin{mdframed}
Recuerde que este es su trabajo, y sólo usted puede expresar con precisión lo que ha aprendido y quiere transmitir. Si lo hace bien, su introducción será más significativa y valiosa que cualquier texto automatizado. ¡Confíe en sus habilidades, y verá que puede hacer un mejor trabajo que cualquier herramienta que automatiza la generación de texto!
\end{mdframed}

